\chapter{The Website Stage}

Where all our efforts goes on the web.  There are two very different
tasks for \FramaC to go online.  Authoring the web site pages is the
responsibility of the developers and the release manager. Publishing
the web site can only be performed by authorized people, who may not
be the release manager.

\section{Requirements}

The website \texttt{README.md} provides a quick guide to locally setup
the website.

\section{Generate website}

Before pushing the branch to the websit, it can be deployed locally but
this is not absolutely necessary.

Push the branch generated by the release script on the website repository.

You can have a look to the generated pages on \texttt{pub.frama-c.com}.
Note however that the download links won't work as they are available only
once the branch is merged (\texttt{download} links are handled by Git-LFS).

Check the following pages:

\begin{itemize}
  \item \texttt{index.html} must display:
  \begin{itemize}
    \item the new event,
    \item a link to the (beta) release at the bottom.
  \end{itemize}
  \item \texttt{/fc-versions/<version\_name>.html}:
  \begin{itemize}
    \item check Changelog link,
    \item check manual links (reminder: the links are dead at this moment), it must contain \texttt{NN.N-Version}
    \item check ACSL version.
  \end{itemize}
  \item \texttt{/html/changelog.html\#Version-NN.N}
  \item \texttt{/html/installations/titanium.html}
  \item \texttt{/html/acsl.html}: check ACSL versions list
  \item \texttt{rss.xml}: check last event
\end{itemize}

On GitLab, the following files must appear as \textbf{new} in \texttt{download}:

\begin{itemize}
  \item \texttt{acsl-X.XX.pdf}
  \item \texttt{acsl-implementation-NN.N-VERSION.pdf}
  \item \texttt{plugin-development-guide-NN.N-VERSION.pdf}
  \item \texttt{user-manual-NN.N-VERSION.pdf}
  \item \texttt{<plugin>-manual-NN.N-VERSION.pdf}\\
        for Aorai, EVA, Metrics, RTE and WP
  \item \texttt{e-acsl/e-acsl-X.XX.pdf}
  \item \texttt{e-acsl/e-acsl-implementation-NN.N-VERSION.pdf}
  \item \texttt{e-acsl/e-acsl-manual-NN.N-VERSION.pdf}
  \item \texttt{aorai-example-NN.N-VERSION.tgz}
  \item \texttt{frama-c-NN.N-VERSION.tar.gz}
  \item \texttt{frama-c-api-NN.N-VERSION.tar.gz}
  \item \texttt{hello-NN.N-VERSION.tar.gz}
\end{itemize}

On GitLab, the following files must appear as \textbf{modified} in \texttt{download}:

\begin{itemize}
  \item \texttt{acsl.pdf}
  \item \texttt{frama-c-acsl-implementation.pdf}
  \item \texttt{frama-c-plugin-development-guide.pdf}
  \item \texttt{frama-c-user-manual.pdf}
  \item \texttt{frama-c-<plugin>-manual.pdf}\\
        for Aorai, EVA, Metrics, RTE, \textbf{Value Analysis} and WP
  \item \texttt{e-acsl/e-acsl.pdf}
  \item \texttt{e-acsl/e-acsl-implementation.pdf}
  \item \texttt{e-acsl/e-acsl-manual.pdf}
  \item \texttt{frama-c-aorai-example.tgz}
  \item \texttt{hello.tar.gz}
\end{itemize}

\section{Announcements}

\begin{itemize}
\item Send an e-mail to \texttt{frama-c-discuss} announcing the release.
\item Tweet the release, pointing to the Downloads page.
\item Ideally, a blog post should arrive in a few days, with some interesting
  new features.
\end{itemize}

\section{Opam package}

You'll need a GitHub account to create a pull request on the official opam repository,
\texttt{ocaml/opam-repository.git}.

\begin{itemize}
\item Clone \texttt{opam-repository}: \texttt{git clone git@github.com:ocaml/opam-repository.git}
\item Make sure you are on \texttt{master} and your branch is up-to-date
\item Create a new directory: \\
  \texttt{packages/frama-c/frama-c.<version>} \\
\item Copy the file \texttt{./distributed/opam} to the opam repository clone in: \\
 \texttt{packages/frama-c/frama-c.<version>/opam}
\item You can provide \verb|sha256| and/or \verb|sha512| checksums at the end of this file if ou wish.
\item (optional) Check locally that everything is fine:
\begin{verbatim}
opam switch create local <some-ocaml-compiler-version>
opam repository add local <path-to-repository-clone>
opam repository set-repos local
opam install frama-c
\end{verbatim}
(of course, if you already create a local switch before and it uses your
local version of the repository, you just have to switch to it).

\textbf{Note:} uncommitted changes are ignored by \texttt{opam repository};
you have to locally commit them before \texttt{opam update} will take them into
account.

\item Create a branch with any name you want (e.g. frama-c.<version>) and push it to your remote Github
\item Create a pull request to opam-repository. If all tests pass,
  someone from opam should merge it for you.
  \textbf{Note:} some opam tests may fail due to external circumstances; if the
  error log makes no sense to you, wait to see if someone will contact you
  about it, or ask directly on the merge request.
\end{itemize}

\section{Updating pub/frama-c}

You'll need to be member of the \texttt{pub/frama-c} project to be able to
commit to it.

\begin{itemize}
\item add the \texttt{pub/frama-c} remote to your Git clone;
\item make sure the last commit is tagged (either a release candidate, or a
  final release);
\item push the stable/\texttt{<codename>} branch to the \texttt{pub} remote.
\end{itemize}

({\em Non-beta only}) After pushing the tag to Gitlab, go to the Releases page
and create a release for the tag.

You should also push the wiki changes. Note that all files listed as \textbf{new}
in the website section should appear in the wiki but:

\begin{itemize}
  \item \texttt{frama-c-api-NN.N-VERSION.tar.gz}
  \item \texttt{hello-NN.N-VERSION.tar.gz}
\end{itemize}

and E-ACSL files are not in a sub-folder \texttt{e-acsl} on the wiki.

\section{Updating the BTS}

\expertise{Usually Julien performs this step. Soon to be obsoleted.}

Each task marked as \texttt{Resolved} and really resolved in this release (say
\texttt{VERSION}) must be marked as \texttt{Fixed in Version VERSION}, then
\texttt{Closed}.

The just released \texttt{VERSION} must be marked as \texttt{released} in the
Section \texttt{manage} while an entry for the next one must be added (if not
already done).

Finally have a look at each still opened task in order to see what should be
resolved in the next release.

%% \section{Update the Wiki}

%% Update references of the Wiki, especially pointers to manuals of this page:
%% \url{http://bts.frama-c.com/dokuwiki/doku.php?id=mantis:frama-c:publications}

\section{Other repositories to update}

Check if other Frama-C (and related) repositories need to be updated:

\begin{itemize}
\item \texttt{acsl-language/acsl} (if last minute patches were applied)
\item \texttt{pub/open-source-case-studies}
\item \texttt{pub/sate-6}
\item other \texttt{pub} repositories related to Frama-C...
\end{itemize}

\section{Preparing the Next Release}

Just update the \texttt{VERSION} file in \texttt{master}, by adding
\texttt{"+dev"}. Do not add any newline at the end of the
\texttt{VERSION} file.

%%% Local Variables:
%%% mode: latex
%%% TeX-master: "release"
%%% End:
